%%% The main file. It contains definitions of basic parameters and includes all other parts.

%% Settings for single-side (simplex) printing
% Margins: left 40mm, right 25mm, top and bottom 25mm
% (but beware, LaTeX adds 1in implicitly)
%\documentclass[12pt,a4paper]{report}
%\setlength\textwidth{145mm}
%\setlength\textheight{247mm}
%\setlength\oddsidemargin{15mm}
%\setlength\evensidemargin{15mm}
%\setlength\topmargin{0mm}
%\setlength\headsep{0mm}
%\setlength\headheight{0mm}
% \openright makes the following text appear on a right-hand page
%\let\openright=\clearpage

%% Settings for two-sided (duplex) printing
\documentclass[12pt,a4paper,twoside,openright]{report}
% \setlength\textwidth{145mm}
% \setlength\textheight{247mm}
% \setlength\oddsidemargin{14.2mm}
% \setlength\evensidemargin{0mm}
% \setlength\topmargin{0mm}
% \setlength\headsep{0mm}
% \setlength\headheight{0mm}
\let\openright=\cleardoublepage

%% Character encoding: usually latin2, cp1250 or utf8:
\usepackage[utf8]{inputenc}

%% Further useful packages (included in most LaTeX distributions)
\usepackage{amsmath}        % extensions for typesetting of math
\usepackage{amsfonts}       % math fonts
\usepackage{amsthm}         % theorems, definitions, etc.
%\usepackage{bbding}         % various symbols (squares, asterisks, scissors, ...)
%\usepackage{bm}             % boldface symbols (\bm)
\usepackage{graphicx}       % embedding of pictures
%\usepackage{fancyvrb}       % improved verbatim environment
%\usepackage[square]{natbib}         % citation style AUTHOR (YEAR), or AUTHOR [NUMBER]
%\usepackage[nottoc]{tocbibind} % makes sure that bibliography and the lists
			    % of figures/tables are included in the table
			    % of contents
%\usepackage{dcolumn}        % improved alignment of table columns
%\usepackage{booktabs}       % improved horizontal lines in tables
%\usepackage{paralist}       % improved enumerate and itemize
\usepackage{tikz}
\usetikzlibrary{shapes,fit,positioning,snakes,mindmap,trees,decorations.text,arrows.meta}
\usepackage{gnuplot-lua-tikz}
\usepackage{nomencl}
\usepackage{epigraph}
\makenomenclature
\usepackage{algorithm,algpseudocode}
%\usepackage[usenames]{xcolor}  % typesetting in color

%%% Basic information on the thesis

% Thesis title in English (exactly as in the formal assignment)
\def\ThesisTitle{Pattern recognition for in-game spell systems}

% Author of the thesis
\def\ThesisAuthor{Pavel Mikuš}

% Year when the thesis is submitted
\def\YearSubmitted{2018}

% Name of the department or institute, where the work was officially assigned
% (according to the Organizational Structure of MFF UK in English,
% or a full name of a department outside MFF)
\def\Department{Department of Software Engineering}

% Is it a department (katedra), or an institute (ústav)?
\def\DeptType{Department}

% Thesis supervisor: name, surname and titles
\def\Supervisor{Miroslav Kratochvíl, M.Sc.}

% Supervisor's department (again according to Organizational structure of MFF)
\def\SupervisorsDepartment{Department of Software Engineering}

% Study programme and specialization
\def\StudyProgramme{Computer Science}
\def\StudyBranch{General Computer Science}

\setlength\epigraphwidth{9cm}
\setlength\epigraphrule{0pt}

% An optional dedication: you can thank whomever you wish (your supervisor,
% consultant, a person who lent the software, etc.)
\def\Dedication{%

\vspace{3 ex}
\noindent
TODO
}

% Abstract (recommended length around 80-200 words; this is not a copy of your thesis assignment!)
\def\Abstract{%
TODO
}

% 3 to 5 keywords (recommended), each enclosed in curly braces
\def\Keywords{%
{spell systems} {shape recognition} {games} {neural network} {shape combination}
}

%% The hyperref package for clickable links in PDF and also for storing
%% metadata to PDF (including the table of contents).
\usepackage[pdftex,unicode]{hyperref}   % Must follow all other packages
\hypersetup{breaklinks=true}
\hypersetup{pdftitle={\ThesisTitle}}
\hypersetup{pdfauthor={\ThesisAuthor}}
\hypersetup{pdfkeywords=\Keywords}
\hypersetup{urlcolor=blue}

% Definitions of macros (see description inside)
%%% This file contains definitions of various useful macros and environments %%%
%%% Please add more macros here instead of cluttering other files with them. %%%

%%% Minor tweaks of style

% These macros employ a little dirty trick to convince LaTeX to typeset
% chapter headings sanely, without lots of empty space above them.
% Feel free to ignore.
\makeatletter
\def\@makechapterhead#1{
  {\parindent \z@ \raggedright \normalfont
   \Huge\bfseries \thechapter. #1
   \par\nobreak
   \vskip 20\p@
}}
\def\@makeschapterhead#1{
  {\parindent \z@ \raggedright \normalfont
   \Huge\bfseries #1
   \par\nobreak
   \vskip 20\p@
}}
\makeatother

% This macro defines a chapter, which is not numbered, but is included
% in the table of contents.
\def\chapwithtoc#1{
\chapter*{#1}
\addcontentsline{toc}{chapter}{#1}
}

% Draw black "slugs" whenever a line overflows, so that we can spot it easily.
\overfullrule=1mm

%%% Macros for definitions, theorems, claims, examples, ... (requires amsthm package)

\theoremstyle{plain}
\newtheorem{thm}{Theorem}
\newtheorem{lemma}[thm]{Lemma}
\newtheorem{claim}[thm]{Claim}

\theoremstyle{plain}
\newtheorem{defn}{Definition}

\theoremstyle{remark}
\newtheorem*{cor}{Corollary}
\newtheorem*{rem}{Remark}
\newtheorem*{observation}{Observation}
\newtheorem*{example}{Example}

%%% An environment for proofs

%%% FIXME %%% \newenvironment{proof}{
%%% FIXME %%%   \par\medskip\noindent
%%% FIXME %%%   \textit{Proof}.
%%% FIXME %%% }{
%%% FIXME %%% \newline
%%% FIXME %%% \rightline{$\square$}  % or \SquareCastShadowBottomRight from bbding package
%%% FIXME %%% }

%%% An environment for typesetting of program code and input/output
%%% of programs. (Requires the fancyvrb package -- fancy verbatim.)

%\DefineVerbatimEnvironment{code}{Verbatim}{fontsize=\small, frame=single}

%%% The field of all real and natural numbers
\newcommand{\R}{\mathbb{R}}
\newcommand{\N}{\mathbb{N}}
\newcommand{\F}{\mathbb{F}}
\newcommand{\Z}{\mathbb{Z}}

%%% Transposition of a vector/matrix
\newcommand{\T}[1]{#1^\top}

%%% Various math goodies
\newcommand{\goto}{\rightarrow}
\newcommand{\gotop}{\stackrel{P}{\longrightarrow}}
\newcommand{\maon}[1]{o(n^{#1})}
\newcommand{\abs}[1]{\left|{#1}\right|}
\newcommand{\dint}{\int_0^\tau\!\!\int_0^\tau}
\newcommand{\isqr}[1]{\frac{1}{\sqrt{#1}}}

%%% Various table goodies
\newcommand{\pulrad}[1]{\raisebox{1.5ex}[0pt]{#1}}
\newcommand{\mc}[1]{\multicolumn{1}{c}{#1}}


% Title page and various mandatory informational pages
\begin{document}
%%% Title page of the thesis and other mandatory pages

%%% Title page of the thesis

\pagestyle{empty}
\hypersetup{pageanchor=false}
\begin{center}

\centerline{\mbox{\includegraphics[width=166mm]{logo-en.pdf}}}

\vspace{-8mm}
\vfill

{\bf\Large BACHELOR THESIS}

\vfill

{\LARGE\ThesisAuthor}

\vspace{15mm}

{\LARGE\bfseries\ThesisTitle}

\vfill

\Department

\vfill

\begin{tabular}{rl}

Supervisor of the bachelor thesis: & \Supervisor \\
\noalign{\vspace{2mm}}
Study programme: & \StudyProgramme \\
\noalign{\vspace{2mm}}
Study branch: & \StudyBranch \\
\end{tabular}

\vfill

% Zde doplňte rok
Prague \YearSubmitted

\end{center}

\newpage

%%% Here should be a bound sheet included -- a signed copy of the "bachelor
%%% thesis assignment". This assignment is NOT a part of the electronic
%%% version of the thesis. DO NOT SCAN.

%%% A page with a solemn declaration to the bachelor thesis

\openright
\hypersetup{pageanchor=true}
\pagestyle{plain}
\pagenumbering{roman}
\vglue 0pt plus 1fill

\noindent
I declare that I carried out this bachelor thesis independently, and only with the cited
sources, literature and other professional sources.

\medskip\noindent
I understand that my work relates to the rights and obligations under the Act No.~121/2000 Sb.,
the Copyright Act, as amended, in particular the fact that the Charles
University has the right to conclude a license agreement on the use of this
work as a school work pursuant to Section 60 subsection 1 of the Copyright Act.

\vspace{10mm}

\hbox{\hbox to 0.5\hsize{%
In .............. date ....................	% FIXME!
\hss}\hbox to 0.5\hsize{%
signature of the author
\hss}}

\vspace{20mm}
\newpage

%%% Mandatory information page of the thesis

\openright

\vbox to 0.5\vsize{
\setlength\parindent{0mm}
\setlength\parskip{5mm}

Title:
\ThesisTitle

Author:
\ThesisAuthor

\DeptType:
\Department

Supervisor:
\Supervisor, \SupervisorsDepartment

Abstract:
\Abstract

Keywords:
\Keywords

\vss}

\newpage

%%% Dedication

\openright

\noindent
\Dedication

\newpage

\openright
\pagestyle{plain}
\pagenumbering{arabic}
\setcounter{page}{1}


%%% A page with automatically generated table of contents of the bachelor thesis

\tableofcontents

%%% Each chapter is kept in a separate file
\chapter{Introduction}
\label{chap:gf}

Magic has been an popular part of computer games since their beginning. In many games, like in any fantasy work, magic has its own lore and laws that make it systematical. Great examples of a complexity of magical spells are games Magicka and Magicka2 //TODO, where player casts spells by combining eight elements. For example, using only earth element results in a rock thrown at the enemy, but adding fire will create a classic fireball.

However, its also usually very easy to do magic in games, since players can cause magical storms by pushing few buttons. This spoils the feeling of magic as something extraordinary and secret. In books and movies we can see wizards performing complicated hand gestures, or drawing on the floor some complex shapes, and magic is restricted to a few chosen, and we would like to offer players something similar.

While majority of games binds spells to buttons and restricts players through artificial systems, like special energy required to cast spells, or cool-down of spell, there already exist several ways of how to make player more connected with the magic. In the game Fable: Adventure //TODO, developers used Kinect technology to track players hands. When player performed a correct gesture, game recognized it and casted corresponding spell. 

Another technique used in several games is to let player draw into a predefined area, or through predefined points. Both ***Castlevania: Dawn of Sorrow ***(some pictures here) and Deep Labyrinth take advantage of Nintendo DS drawing interface that allows players draw magic signs. In Castlevania, players draw signs by connecting glowing points in circle in out of combat situations. Deep Labyrinth introduces special casting interface as well, consisting of 3x3 sized grid, where player connects dots. These approaches take away part of the freedom, but they make recognition algorithms much easier, e.g. by tracing only the order of points passed and selecting the right spell.

Different approach to recognition of player drawn spells and their recognition was used in Arx Fatalis. Players were drawing symbols using their mouse, and the sequence of symbols represented some spell. While casting, game encoded mouse moves in 8-directions precision, each direction representing some letter. After player finished their spell, a Levenshtein distance was calculated from each spell to the user created sequence, and best candidate (lowest Levenshtein cost) was returned. 

\chapter{Goals}
\label{chap:gf}

We would like to allow players draw complex spells. Instead of time sequence of symbols, players could draw symbols arranged in a shape of other symbol, or embed one symbol in another. We also want our approach to be applicable in multi-player environments, where new situations can occur, such as multiple players drawing one spell. We don't want to introduce aiding structures such as trigger points, since they require order in which they were passed and force player to draw in a certain way or direction. They are also not very usable in multi-player environment, since they can't handle multiple people drawing at the same time.

For the purpose of this work, we specify following requirements that we have on our recognition system. 

Durability against shape deformations: Since we want to recognize hand-drawn shapes, our system has to be prepared against human imprecise drawing, especially when drawing with mouse.

Extendability: We would like to offer easy-to-use library, that other developers can incorporate. For this purposes we would like to let them define their own shapes that should be recognized.

Speed: Our system should be applicable in demanding video games environment, prepared for the possibility of many players drawing their spells. To achieve this, we require fast recognition technique as well as usability in parallel environment.

Recognition of embedded shapes and shape conglomerations: To allow players cast complex spells, we need to give them ability to somehow encode multiple symbols in one spell. We have decided that it will be in a form of shape embeddings, which means that shape can contain another shape in defined spot. We also allow them to create shape conglomerations, where several shapes are organized in a shape of symbol. TODO examples

For these purposes, we need to find pattern recognition algorithm that most closely fits our needs and modify it to eliminate any of its shortcomings.

\chapter{Pattern Recognition}
\label{chap:gf}

There are thousands of algorithms for pattern recognition as it is central problem for . We can divide them into learning-based approaches and template-based approaches [TODO skeletons.pdf p.193].in the learning-based approaches, pattern classifiers are obtained trough training using classified samples. In the template based approaches, patterns are described by templates and the recognition problem becomes searching for the best matching template for a given input image.

[TODO distanceTransform.pdf p.1]
We can also divide them based on the level of image feature extraction into three classes: algorithms that use pixel values directly, e.g. correlation methods; algorithms that use low level features such as edges and corners, e.g. distance transform method; and algorithms that use high level features such as identified objects or relation between the features, e.g. graph-theoretic-methods. 

To make this work possible, we have selected only several algorithms from different categories that we briefly explain and compare with our requirements.

\subchapter{Template matching}




\include{epilog}

%%% Bibliography
%%% Bibliography (literature used as a source)
%%%
%%% We employ bibTeX to construct the bibliography. It processes
%%% citations in the text (e.g., the \cite{...} macro) and looks up
%%% relevant entries in the bibliography.bib file.
%%%
%%% The \bibliographystyle command selects, which style will be used
%%% for references from the text. The argument in curly brackets is
%%% the name of the corresponding style file (*.bst). Both styles
%%% mentioned in this template are included in LaTeX distributions.

% \bibliographystyle{plainnat}    %% Author (year)
% \bibliographystyle{unsrt}     %% [number]
\bibliographystyle{alpha}

\renewcommand{\bibname}{Bibliography}

%%% Generate the bibliography. Beware that if you cited no works,
%%% the empty list will be omitted completely.

\bibliography{bibliography}

%%% If case you prefer to write the bibliography manually (without bibTeX),
%%% you can use the following. Please follow the ISO 690 standard and
%%% citation conventions of your field of research.

% \begin{thebibliography}{99}
%
% \bibitem{lamport94}
%   {\sc Lamport,} Leslie.
%   \emph{\LaTeX: A Document Preparation System}.
%   2nd edition.
%   Massachusetts: Addison Wesley, 1994.
%   ISBN 0-201-52983-1.
%
% \end{thebibliography}


\listofalgorithms

%%% Figures used in the thesis (consider if this is needed)
\listoffigures

%%% Tables used in the thesis (consider if this is needed)
%%% In mathematical theses, it could be better to move the list of tables to the beginning of the thesis.
\listoftables

%%% Abbreviations used in the thesis, if any, including their explanation
%%% In mathematical theses, it could be better to move the list of abbreviations to the beginning of the thesis.
%\chapwithtoc{List of Abbreviations}
\renewcommand{\nomname}{List of Abbreviations}
%\printnomenclature

%%% Attachments to the bachelor thesis, if any. Each attachment must be
%%% referred to at least once from the text of the thesis. Attachments
%%% are numbered.
%%%
%%% The printed version should preferably contain attachments, which can be
%%% read (additional tables and charts, supplementary text, examples of
%%% program output, etc.). The electronic version is more suited for attachments
%%% which will likely be used in an electronic form rather than read (program
%%% source code, data files, interactive charts, etc.). Electronic attachments
%%% should be uploaded to SIS and optionally also included in the thesis on a~CD/DVD.
\chapwithtoc{Attachments}

\section*{Attachment A - the Enclosed CD}

On the CD attached to this thesis (and on the online Thesis Repository of Charles University \cite{TRepo}) we enclose
the source codes of the implemented software together with the source codes of its
dependencies TODO

The electronic version of this thesis is also enclosed.

\openright
\end{document}
